\documentclass[12pt,a4paper]{article}
\usepackage[utf8]{inputenc}
\usepackage[T1]{fontenc}
\usepackage[brazil]{babel}
\usepackage{geometry}
\usepackage{hyperref}
\usepackage{graphicx}
\usepackage{listings}
\usepackage{xcolor}
\usepackage{booktabs}
\usepackage{tabularx}
\usepackage{longtable}
\usepackage{float}
\geometry{margin=2.5cm}

\hypersetup{
  colorlinks=true,
  linkcolor=blue,
  urlcolor=blue,
  citecolor=blue
}

\lstdefinestyle{code}{
  basicstyle=\ttfamily\small,
  keywordstyle=\color{blue!70!black},
  commentstyle=\color{green!50!black},
  stringstyle=\color{red!60!black},
  showstringspaces=false,
  breaklines=true,
  frame=single,
  framerule=0.3pt,
  numbers=left,
  numberstyle=\tiny, 
  xleftmargin=1em,
}

\title{Arquitetura e Implementação do Backend de um Jogo Musical Web: \textit{web2-musicalGame}}
\author{Autor: Equipe do Projeto}
\date{Setembro de 2025}

\begin{document}
\maketitle

\begin{abstract}
Este artigo descreve em detalhe a arquitetura, o desenho das APIs e as principais decisões de implementação do backend do projeto \textit{web2-musicalGame}, bem como sua comunicação com o frontend React. O sistema provê autenticação por JWT, gerenciamento de desafios musicais com integração a metadados externos (Deezer/Spotify), mecânica de jogo com avaliação de respostas (\textit{fuzzy matching}) e persistência de tentativas apenas ao final de cada execução. Também discutimos sanitização/validação de entradas, agregações para ranking e perfis, e diretrizes para execução local e implantação.
\end{abstract}

\section{Introdução}
Aplicações educacionais e de entretenimento musical se beneficiam de mecânicas de adivinhação com trechos de áudio curtos e progressão por dificuldade. O \textit{web2-musicalGame} oferece uma experiência de jogo onde usuários tentam identificar músicas a partir de prévias de áudio, com pontuação dinâmica e registro de desempenho. Este artigo foca no backend (Node.js + Express + Prisma) e na forma como expõe serviços consumidos pelo frontend (Vite + React).

\section{Trabalhos Relacionados}
Jogos de \textit{quiz} musical e plataformas de prática auditiva frequentemente utilizam APIs de metadados (p.~ex., Deezer/Spotify) e regras de comparação aproximada de texto para aceitar variações de grafia. Este projeto incorpora tais práticas, adicionando uma estratégia de persistência \textit{após conclusão} do desafio que simplifica consistência de ranking.

\section{Arquitetura Geral}
\subsection{Visão de alto nível}
O backend é um serviço REST em Node.js (Express), com persistência por Prisma ORM e banco SQLite ou PostgreSQL. A autenticação usa JWT e o frontend consome os endpoints via Axios.

\begin{itemize}
  \item \textbf{Camada de Roteamento e Controladores}: mapeia endpoints para casos de uso (\texttt{auth}, \texttt{desafio}, \texttt{desafioMusica}, \texttt{game}, \texttt{ranking}, \texttt{profile}, \texttt{deezer}, \texttt{spotify}, \texttt{meta}).
  \item \textbf{Middlewares}: autenticação e autorização (\textit{roles}).
  \item \textbf{Serviços}: integração com provedores externos (p.~ex., Deezer) e utilidades de domínio.
  \item \textbf{Utilitários}: validação e sanitização (limites de tamanho, normalização de dificuldade, parsing seguro de inteiros).
\end{itemize}

\subsection{Modelagem de dados (Prisma)}
A seguir, um trecho do esquema (\texttt{schema.prisma}) com entidades centrais:\footnote{No repositório, o provider pode ser \texttt{postgresql} ou \texttt{sqlite}; a URL do banco é configurada via \texttt{DATABASE\_URL}.}

\begin{lstlisting}[style=code, language=]
model Usuario {
  id         Int         @id @default(autoincrement())
  nome       String
  email      String      @unique
  senha      String?
  tipo       String      // "admin" ou "usuario"
  provider   String      @default("local")
  picture    String      @default("https://cdn-icons-png.flaticon.com/512/17/17004.png")
  tentativas Tentativa[]
}

model Desafio {
  id          Int              @id @default(autoincrement())
  titulo      String
  genero      String
  dificuldade Dificuldade
  tentativas  Tentativa[]
  musicas     DesafioMusica[]
  desafioCapa String?
  createdAt   DateTime         @default(now())
  updatedAt   DateTime         @updatedAt
}

model Tentativa {
  id            Int      @id @default(autoincrement())
  usuarioId     Int
  desafioId     Int
  acertou       Boolean
  tempoResposta Float?
  pontos        Int
  usuario       Usuario  @relation(fields: [usuarioId], references: [id], onDelete: Cascade)
  desafio       Desafio  @relation(fields: [desafioId], references: [id], onDelete: Cascade)
}

model DesafioMusica {
  id         Int     @id @default(autoincrement())
  desafioId  Int
  deezerId   String
  desafio    Desafio @relation(fields: [desafioId], references: [id], onDelete: Cascade)
  createdAt  DateTime @default(now())

  @@unique([desafioId, deezerId])
}

enum Dificuldade {
  FACIL
  MEDIO
  DIFICIL
  MUITO_DIFICIL
  EXTREMO
}
\end{lstlisting}

\section{Endpoints e Contratos}
\subsection{Autenticação}
\begin{itemize}
  \item POST \texttt{/auth/register}, \texttt{/auth/login}, \texttt{/auth/google}: retorna \texttt{\{ user, token \}}
  \item Middleware: extrai JWT (\texttt{Authorization: Bearer <token>}) e injeta \texttt{req.usuario}; rotas admin validam \texttt{usuario.tipo}.
\end{itemize}

\subsection{Administração de Desafios}
\begin{itemize}
  \item GET \texttt{/desafios}: lista com \texttt{musicasCount}
  \item POST/PUT/DELETE \texttt{/desafios[/:id]}: criação/edição/exclusão com validação e sanitização
  \item POST/DELETE \texttt{/desafio-musica}: vincula ou remove música, retornando \texttt{musicas} e \texttt{musicasCount}
  \item GET \texttt{/deezer/search}, GET \texttt{/spotify/search}: sugestões para adicionar faixas
\end{itemize}

\subsection{Mecânica de Jogo}
\begin{itemize}
  \item POST \texttt{/game/start/:desafioId}: inicia sessão de jogo em memória e retorna primeira rodada
  \item GET \texttt{/game/current/:sessionId}: retorna estado atual
  \item POST \texttt{/game/guess/:sessionId}: avalia palpite, pontua e avança; ao fim, persiste tentativas
\end{itemize}

\subsection{Ranking e Perfil}
\begin{itemize}
  \item GET \texttt{/ranking/global}: agregação de pontos por usuário considerando somente execuções completas
  \item GET \texttt{/profile/me}: métricas individuais (taxa de acerto, média de tempo, distribuição por dificuldade)
\end{itemize}

\section{Regras de Negócio e Implementação}
\subsection{Validação e Sanitização}
Entradas passam por \textit{sanitizers}: limite de tamanho (e remoção de caracteres de controle), \textit{normalização de dificuldade}, \textit{parse} seguro de inteiros e limpeza de consultas simples (busca Deezer/Spotify). URLs de capa aceitam http/https/data/blob/ipfs e caminhos relativos.

\subsection{Jogo e Pontuação}
A sessão de jogo mantém estado em memória: música atual, contadores de \textit{replay}, tempo decorrido, pontuação acumulada e \textit{streak}. A validação de acerto usa normalização (acentos, pontuação, parênteses) e \textit{fuzzy matching} (substrings, tokens e distância de Levenshtein com limiares por tamanho). Pontos combinam base por dificuldade e bônus por rapidez.

Ao finalizar o desafio, o backend persiste \textbf{todas as tentativas} daquela execução (uma por música). Se o usuário desistir antes do fim, nada é gravado, simplificando consistência para ranking.

\subsection{Ranking}
A agregação considera apenas execuções completas. Pontos por desafio/usuário são contabilizados uma única vez (primeira execução completa), enquanto métricas descritivas (acertos, acurácia, média de tempo) são derivadas das tentativas persistidas.

\section{Comunicação com o Frontend}
\subsection{Padrões de Resposta}
O backend retorna DTOs estáveis e erros padronizados (\texttt{code}, \texttt{message}). O frontend usa Axios com baseURL e interceptador 401 para \textit{logout}. Ao adicionar/remover músicas de um desafio, o backend já retorna lista atualizada e \texttt{musicasCount}, evitando \textit{overfetch}.

\subsection{Fluxos Exemplares}
\paragraph{Adicionar música no admin} Busca com debounce em \texttt{/deezer/search}; seleção dispara \texttt{/desafio-musica} e a resposta atualiza a UI (lista + contador).

\paragraph{Jogar um desafio} Lista de desafios → \texttt{/game/start} → rodadas com avaliação em \texttt{/game/guess} → persistência ao fim → tela de resumo.

\section{Execução e Implantação}
\subsection{Ambiente de Desenvolvimento}
\begin{enumerate}
  \item Configurar \texttt{backend/.env} com \texttt{DATABASE\_URL}, \texttt{JWT\_SECRET} e porta
  \item Rodar migrações: \texttt{npx prisma migrate dev} e gerar client: \texttt{npx prisma generate}
  \item Iniciar backend (\texttt{npm run dev} em \texttt{backend/}) e frontend (\texttt{npm run dev} na raiz)
\end{enumerate}

\subsection{Banco de Dados}
Suporte a SQLite (\texttt{file:./dev.db}) e PostgreSQL (\texttt{postgresql://...}). Em mudanças estruturais, usar \texttt{migrate dev} ou \texttt{migrate reset} em ambiente local.

\section{Discussão}
As decisões de persistir somente após a conclusão do desafio, empacotar contadores no DTO e consolidar a validação/sanitização no backend reduzem complexidade no cliente e previnem inconsistências. O uso de Prisma oferece produtividade e segurança de tipos, enquanto o desenho dos endpoints favorece uma experiência fluida no frontend (UI otimista e sobreposições de feedback unificadas).

\section{Conclusão}
Apresentamos a arquitetura e implementação do backend do \textit{web2-musicalGame}, detalhando suas camadas, modelos, regras e contratos com o frontend. Como trabalhos futuros, sugerimos adicionar testes automatizados (unitários e de integração), métricas de observabilidade e suporte a sessões distribuídas.

\section*{Agradecimentos}
Agradecemos à equipe e aos docentes envolvidos no desenvolvimento e testes do projeto.

\section{Camadas do Backend com Exemplos}
\subsection{Servidor Express e Rotas}
\begin{lstlisting}[style=code, language=JavaScript]
// backend/src/server.js
import express from "express";
import cors from "cors";
import process from "process";
import authRoutes from "./routes/auth.routes.js";
import desafioRoutes from "./routes/desafio.routes.js";
import spotifyRoutes from "./routes/spotify.routes.js";
import desafioMusicaRoutes from "./routes/desafioMusica.routes.js";
import metaRoutes from "./routes/meta.routes.js";
import gameRoutes from "./routes/game.routes.js";
import deezerRoutes from './routes/deezer.routes.js';
import profileRoutes from './routes/profile.routes.js';
import rankingRoutes from './routes/ranking.routes.js';

const app = express();
app.use(cors());
app.use(express.json());
app.use("/auth", authRoutes);
app.use("/desafios", desafioRoutes);
app.use("/spotify", spotifyRoutes);
app.use("/desafio-musica", desafioMusicaRoutes);
app.use("/meta", metaRoutes);
app.use('/game', gameRoutes);
app.use('/deezer', deezerRoutes);
app.use('/profile', profileRoutes);
app.use('/ranking', rankingRoutes);
app.get("/ping", (req, res) => res.json({ message: "Servidor rodando!" }));
app.listen(process.env.PORT||3000);
\end{lstlisting}

\subsection{Middleware de Autenticação (JWT)}
\begin{lstlisting}[style=code, language=JavaScript]
// backend/src/middlewares/auth.middleware.js
import { verificarToken } from "../utils/token.js";
export const autenticar = async (req, res, next) => {
  const authHeader = req.headers.authorization;
  if (!authHeader) return res.status(401).json({ erro: "Token não fornecido" });
  const [tipo, token] = authHeader.split(" ");
  if (tipo !== "Bearer" || !token) return res.status(401).json({ erro: "Token mal formatado" });
  try { req.usuario = verificarToken(token); next(); }
  catch { return res.status(401).json({ erro: "Token inválido" }); }
}
\end{lstlisting}

\subsection{Validação e Sanitização}
\begin{lstlisting}[style=code, language=JavaScript]
// backend/src/utils/validation.js
export const LIMITES = { MAX_TITULO: 100, MAX_GENERO: 50, MAX_MUSICAS_DESAFIO: 50 };
export function limitarTamanho(v, max=255) {
  if (typeof v !== 'string') return '';
  let s=''; for (let i=0;i<v.length;i++){ const c=v.charCodeAt(i); if (c<32 && c!==9 && c!==10 && c!==13) continue; s+=v[i]; }
  s=s.trim(); return s.length>max? s.slice(0,max): s;
}
export function sanitizeBuscaSimples(v,max=80){ if(typeof v!=='string') return ''; let s=v.normalize('NFD').replace(/[^\p{L}\p{N}\s'\-_.]/gu,'').replace(/\s+/g,' ').trim(); return s.length>max? s.slice(0,max): s; }
export function normalizarDificuldade(v){ const up=String(v||'').trim().toUpperCase(); const ok=['FACIL','MEDIO','DIFICIL','MUITO_DIFICIL','EXTREMO']; return ok.includes(up)? up:'FACIL'; }
\end{lstlisting}

\subsection{Controller de Desafios}
\begin{lstlisting}[style=code, language=JavaScript]
// backend/src/controllers/desafio.controller.js (trecho)
function toDesafioDTO(d){ return { id:d.id, titulo:d.titulo, genero:d.genero, dificuldade:d.dificuldade, desafioCapa:d.desafioCapa, musicasCount:d._count?.musicas??0, createdAt:d.createdAt, updatedAt:d.updatedAt }; }
export const listarDesafios = async (req,res)=>{ const ds= await prisma.desafio.findMany({ include:{ _count:{ select:{ musicas:true } } }, orderBy:{ id:'desc' } }); res.json(ds.map(toDesafioDTO)); };
export const criarDesafio = async (req,res)=>{ let { titulo, genero, dificuldade, desafioCapa, capa }=req.body; if(!desafioCapa&&capa) desafioCapa=capa; titulo=sanitizeTexto(titulo,LIMITES.MAX_TITULO); genero=sanitizeTexto(genero,LIMITES.MAX_GENERO); dificuldade=normalizarDificuldade(dificuldade); /* valida capa... */ const novo= await prisma.desafio.create({ data:{ titulo, genero, dificuldade, desafioCapa: desafioCapa||null }, include:{ _count:{ select:{ musicas:true } } } }); res.status(201).json(toDesafioDTO(novo)); };
\end{lstlisting}

\subsection{Controller de Jogo e Regras}
\begin{lstlisting}[style=code, language=JavaScript]
// backend/src/controllers/game.controller.js (trecho)
const DIFFICULTY_CONFIG = { FACIL:{snippetSeconds:10,base:8,maxBonus:12}, MEDIO:{snippetSeconds:7,base:12,maxBonus:18}, DIFICIL:{snippetSeconds:5,base:18,maxBonus:27}, MUITO_DIFICIL:{snippetSeconds:3,base:26,maxBonus:39}, EXTREMO:{snippetSeconds:2,base:38,maxBonus:57} };
function normalizarTexto(t){ return (t||'').toLowerCase().normalize('NFD').replace(/[^\p{L}\p{N}\s]/gu,'').replace(/\s+/g,' ').trim(); }
export async function submitGuess(req,res){ const { sessionId }=req.params; const answerSan=limitarTamanho(req.body.answer,120); const s=sessions.get(sessionId); const track=s.tracks[s.current]; const tituloBase=track.titulo.replace(/\([^)]*\)/g,' ');
  const a=normalizarTexto(answerSan), b=normalizarTexto(tituloBase); let correta=false; if(a===b||b.includes(a)||a.includes(b)) correta=true; else { /* tokens + levenshtein + similarity */ }
  // calcula bonus por tempo (janela = snippetSeconds), soma a score e acumula tentativa; persiste no final da execução.
  res.json({ correta, titulo: track.titulo, artista: track.artista, score: s.score, round: s.current+1, totalRounds: s.tracks.length, finished: s.current>=s.tracks.length }); }
\end{lstlisting}

\section{Comunicação Frontend↔Backend com Exemplos}
\subsection{Cliente HTTP}
\begin{lstlisting}[style=code, language=JavaScript]
// src/services/api.js
import axios from 'axios';
export const api = axios.create({ baseURL: 'http://localhost:3000' });
// AuthProvider injeta Authorization quando disponível
\end{lstlisting}

\subsection{Hook de Administração de Desafios}
\begin{lstlisting}[style=code, language=JavaScript]
// src/hooks/useDesafios.js (trecho)
export function useDesafios(){
  const [desafios,setDesafios]=useState([]);
  const carregar=useCallback(async()=>{ const {data}=await api.get('/desafios'); setDesafios(data); },[]);
  const criar=useCallback(async(p)=>{ const {data}=await api.post('/desafios',p); setDesafios(prev=>[data,...prev]); return data; },[]);
  const atualizar=useCallback(async(id,p)=>{ const {data}=await api.put(`/desafios/${id}`,p); setDesafios(prev=> prev.map(d=> d.id===id? {...d,...data}: d)); return data; },[]);
  const remover=useCallback(async(id)=>{ await api.delete(`/desafios/${id}`); setDesafios(prev=> prev.filter(d=> d.id!==id)); },[]);
  return { desafios, carregar, criar, atualizar, remover };
}
\end{lstlisting}

\subsection{Fluxo de Jogo no Frontend}
\begin{lstlisting}[style=code, language=JavaScript]
// src/pages/PlayDesafio.jsx (trecho)
useEffect(()=>{ (async()=>{ const {data}= await api.post(`/game/start/${id}`); setSession(data); })(); },[id]);
async function submit(e){ e.preventDefault(); const { data } = await api.post(`/game/guess/${session.sessionId}`, { answer }); setFeedback({ correta:data.correta, titulo:data.titulo, artista:data.artista }); if(data.finished){ /* exibir resumo */ } else { const { data: next } = await api.get(`/game/current/${session.sessionId}`); setCurrent(next); } }
\end{lstlisting}

\section{Contratos e Erros Padronizados}
\begin{itemize}
  \item Sucesso: objetos DTO estáveis (\texttt{DesafioDTO}, listas de músicas com \texttt{musicasCount}).
  \item Erro: \texttt{\{ error: { code, message } \}} com HTTP adequado (400/401/403/404/500).
\end{itemize}

\section{Execução Local}
\begin{lstlisting}[style=code, language=]
# backend/.env (SQLite)
DATABASE_URL="file:./dev.db"
JWT_SECRET="sua_chave"
PORT=3000

# Migrações e client
npx prisma migrate dev
npx prisma generate

# Rodar
npm run dev  # em backend/
# e na raiz
npm run dev
\end{lstlisting}

\section{Tabela de Endpoints da API}
\begingroup
\small
\setlength{\tabcolsep}{4pt}% menos espaço horizontal entre colunas
\renewcommand{\arraystretch}{1.15}% linhas um pouco mais altas
\begin{longtable}{@{} l l c p{0.32\textwidth} p{0.40\textwidth} @{}}
\caption{Resumo dos principais endpoints, com contratos de requisição e resposta.}\label{tab:endpoints}\\
\toprule
\textbf{Método} & \textbf{Rota} & \textbf{Auth} & \textbf{Corpo (req)} & \textbf{Resposta (res)} \\
\midrule
\endfirsthead
\toprule
\textbf{Método} & \textbf{Rota} & \textbf{Auth} & \textbf{Corpo (req)} & \textbf{Resposta (res)} \\
\midrule
\endhead
\bottomrule
\endfoot
\bottomrule
\endlastfoot

POST & /auth/register & -- & \texttt{\{ nome, email, senha \}} & \texttt{\{ user, token \}} \\
POST & /auth/login & -- & \texttt{\{ email, senha \}} & \texttt{\{ user, token \}} \\
POST & /auth/google & -- & \texttt{\{ credential \}} & \texttt{\{ user, token \}} \\
\midrule
GET & /desafios & opcional & -- & \texttt{[ DesafioDTO ]} \\
POST & /desafios & admin & \texttt{\{ titulo, genero, dificuldade, desafioCapa? \}} & \texttt{DesafioDTO} \\
PUT & /desafios/:id & admin & \texttt{\{ titulo?, genero?, dificuldade?, desafioCapa? \}} & \texttt{DesafioDTO} \\
DELETE & /desafios/:id & admin & -- & \texttt{\{ message, desafio \}} \\
\midrule
GET & /desafio-musica/:id & admin & -- & \texttt{[ MusicaDTO ]} \\
POST & /desafio-musica/:id & admin & \texttt{\{ deezerId \}|\{ musicaNome, artistaNome? \}} & \texttt{\{ musicas: [MusicaDTO], musicasCount \}} \\
DELETE & /desafio-musica/:id & admin & \texttt{\{ deezerId \}|\{ musicaNome, artistaNome? \}} & \texttt{\{ ok, musicasCount, musicas: [MusicaDTO] \}} \\
\midrule
GET & /deezer/search & admin & \texttt{?nome, ?artista} & \texttt{[ MusicaDTO ]} \\
GET & /spotify/search & admin & \texttt{?query} & \texttt{[ ... ]} \\
\midrule
POST & /game/start/:desafioId & user & -- & \texttt{\{ sessionId, totalRounds, dificuldade, snippetSeconds \}} \\
GET & /game/current/:sessionId & user & -- & \texttt{\{ round, totalRounds, preview, artista, imagem, snippetSeconds, score \}} \\
POST & /game/guess/:sessionId & user & \texttt{\{ answer \}} & \texttt{\{ correta, matchType?, titulo, artista, score, pontosGanhos?, pontosBase?, bonusTempo?, round, totalRounds, tempoResposta, finished? \}} \\
\midrule
GET & /ranking/global & -- & \texttt{?limit, ?offset, ?dificuldade, ?periodo} & \texttt{[ \{ usuarioId, nome, pontos, acertos, mediaTempo, acuracia \} ]} \\
GET & /profile/me & user & -- & \texttt{\{ ...stats \}} \\
\end{longtable}
\endgroup

\section{Diagrama de Sequ\^encia (Fluxo de Jogo)}
Para fins de compila\c{c}\~ao simples no Overleaf sem pacotes extras, representamos o diagrama de sequ\^encia como pseudoc\'odigo.

\begin{lstlisting}[style=code, language=]
Participantes: Frontend, API, Prisma/DB, Deezer

Frontend -> API: POST /game/start/:desafioId
API -> Prisma/DB: findUnique(desafio, include: musicas)
loop para cada musica
  API -> Deezer: buscarMusicaDeezerPorId(deezerId)
end
API: embaralha(tracks), cria sessão em memória
API -> Frontend: { sessionId, totalRounds, dificuldade, snippetSeconds }

Frontend -> API: GET /game/current/:sessionId
API: recupera sessão, marca roundStartedAt
API -> Frontend: { round, totalRounds, preview, artista, imagem, snippetSeconds, score }

Frontend -> API: POST /game/guess/:sessionId { answer }
API: normaliza(answer), compara(titulo)
alt correta
  API: calcula bonusTempo e pontosGanhos; score +=
else erro
  API: pontosGanhos = 0
end
API: acumula tentativa em memória (guesses)
API: avança s.current; se acabou -> finished=true e persiste tentativas createMany
API -> Prisma/DB: createMany(tentativas[]) (somente ao final)
API -> Frontend: { correta, matchType?, titulo, artista, score, round, totalRounds, tempoResposta, finished? }

Frontend: se finished -> exibe resumo e botões Reiniciar/Outro desafio/Home
\end{lstlisting}

\end{document}
